%=============================================================================
\section{Background and Related Work}
\label{sec:related}
%=============================================================================

\subsection{Searching with Errors}

Identifying an unknown element through queries whose answers may be wrong
originates with R\'enyi's guessing game in which the responder lies with fixed
probability~\cite{renyi1961}, and independently with Ulam's question of how
many yes/no queries suffice to locate an integer below one million when a
bounded number of answers may be false~\cite{ulam1976}. The combinatorial
branch, in which lies are adversarial and bounded in number, produces optimal
strategies with a direct correspondence to error-correcting codes; Pelc's
survey remains the standard reference~\cite{pelc2002}. The probabilistic
branch treats errors as channel noise and connects to sequential transmission
with feedback~\cite{horstein1963,burnashev1974}, to noisy binary
search~\cite{karp2007,nowak2011}, and to treatments in which the noise depends
on the query itself~\cite{kaspi2018,chung2018}.

Our setting sits in the probabilistic branch, with one departure that turns
out to matter. Existing formulations assume a single crossover probability, or
make it depend on a continuous query parameter. In an attribute catalogue the
crossover is a property of the semantic category being probed and can vary by
an order of magnitude across categories. That heterogeneity, rather than the
presence of noise as such, is what drives the effect we report.

\subsection{Active Learning and Experimental Design under Noise}

Choosing the experiment that maximises expected information about a latent
quantity is the central prescription of Bayesian experimental
design~\cite{lindley1956,chaloner1995}, and in the noiseless binary case it
reduces to maximising the entropy of the induced partition. Geman and
Jedynak's active testing model~\cite{geman1996active} is an early application
of the same principle to sequential narrowing. Jedynak et
al.~\cite{jedynak2012} give the result most often invoked here: for twenty
questions with noise under entropy loss, a greedy one-step policy is optimal
over the whole horizon.

That result is frequently cited as licence to select greedily and stop
worrying, and it is worth being precise about what it does and does not cover.
It is established under a stated observation model with a single noise level;
it is not a statement about the heterogeneous, per-attribute channel studied
here. The active-learning literature is in fact explicit that noise breaks
naive greedy selection: Golovin, Krause and Ray show that generalised binary
search, near-optimal without noise, ``can perform very poorly'' once
observations are noisy, and give EC$^2$, a greedy algorithm over equivalence
classes with a competitiveness guarantee under
noise~\cite{golovin2010noisy}; the adaptive-submodularity framework supplies
the general machinery~\cite{golovin2011adaptive}. Our position relative to
this work should be stated plainly, since it is the closest technical prior
art. EC$^2$ and its relatives take the noise model as given and construct a
selection rule with guarantees against it. We take the selection rule as given
--- greedy mutual information, which is what deployed systems actually use ---
and make an empirical claim about the \emph{shape} of the noise model in real
attribute catalogues, namely that it is heterogeneous across attributes,
partly manufactured by preprocessing, and aligned against the naive
criterion. The two are complementary: a practitioner adopting EC$^2$ would
still need to know the per-attribute crossover probabilities, and
Section~\ref{sec:reliability} is about how to get them.

The broader active-learning
literature~\cite{settles2009active} treats oracle fallibility under the
heading of proactive learning, where oracles differ in cost and
reliability and the learner chooses among them~\cite{donmez2008proactive}, and
under label noise more generally~\cite{frenay2014label}. The structural
difference in our setting is that there is one oracle whose reliability varies
by \emph{question}, not several oracles of differing quality.

\subsection{Estimating Reliability from Repeated Observations}

The problem of recovering an annotator's error rate without ground truth is
old and well solved. Dawid and Skene's EM procedure estimates per-observer
confusion matrices jointly with the latent true labels~\cite{dawid1979}, and
the crowdsourcing literature extends it to weighted aggregation and to
labeller-quality models~\cite{sheng2008,whitehill2009}. This matters here for
two reasons. First, it is the route by which we obtain the parameter our method
consumes without assuming it. A deployed system accumulating interaction logs is
in exactly the repeated-observation regime those methods assume, with one
substitution: the latent label is the \emph{item the user has in mind} rather
than the class of an annotated example, and the ``annotators'' are the questions
about each attribute. Section~\ref{sec:reliability-em} makes that substitution
and Section~\ref{sec:results-em} reports that the resulting estimate is good
enough to replace the annotation outright. Second, it frames what our
repeated-item estimator of
Section~\ref{sec:reliability-empirical} is and is not: it applies the same
logic---two independent observations of one underlying entity bound the
instability of the labelling---to the \emph{catalogue} rather than to users,
and therefore measures a lower bound on the quantity of interest rather than
the quantity itself.

\subsection{Interactive and Conversational Recommendation}

Systems that alternate between asking about attributes and proposing items
form a substantial literature, surveyed
in~\cite{gao2021advances,jannach2021survey}, including policy learning over
knowledge graphs~\cite{lei2020ear,lei2020cpr,deng2021unicorn}.
Maximum-entropy attribute querying appears throughout as the standard baseline
against which learned policies are compared. This positions our work
carefully: we are not proposing a competitor to learned conversational
policies but a correction to the baseline those policies are measured against,
which is a modest claim and we make no larger one. The correction is
orthogonal and could be applied inside a learned policy's action-scoring step.
Because the claim that such policies inherit the bias is an empirical one
about a family of systems, Section~\ref{sec:baselines} implements and trains a
member of that family rather than citing it.

\subsection{Question Asking in Language Models}

Recent work evaluates and improves the information-seeking behaviour of
language models, using twenty-questions style tasks as the
testbed~\cite{debruyn2022,mazzaccara2024,hu2024uot}. These systems generate
questions in open natural language rather than selecting from a fixed
attribute-value grid, which makes them more flexible and considerably more
expensive per turn. The bias we describe should apply to them as well, since
expected-information-gain reranking of generated
questions~\cite{mazzaccara2024} inherits the same blind spot, but verifying
that is beyond the present scope and we flag it as future work rather than
claiming it.

\subsection{Positioning}

Table~\ref{tab:related} places this work against the closest prior art. The
honest summary is that the acquisition function we derive is standard mutual
information under a stated observation model, and its derivation is
unsurprising to anyone working in Bayesian experimental design. The
contribution we claim is empirical and diagnostic: that the observation model
in some attribute catalogues is heterogeneous in a specific and unfavourable
way, that this is partly an artefact of preprocessing rather than of the
domain, that it does not hold everywhere and we say where it fails, and that
correcting it requires less knowledge than one might expect.

\begin{table}[t]
\caption{Positioning against the closest prior work. ``Heterog.\ $\varepsilon$''
asks whether the noise level is allowed to differ per attribute.}
\label{tab:related}
\centering
\tiny
\begin{tabular}{@{}lccc@{}}
\toprule
\textbf{Work} & \textbf{Noise} & \textbf{Heterog.\ $\varepsilon$} & \textbf{Selection} \\
\midrule
R\'enyi~\cite{renyi1961}, Ulam~\cite{ulam1976} & yes & no & theory \\
Pelc~\cite{pelc2002} (survey) & yes & no & optimal \\
Jedynak et al.~\cite{jedynak2012} & yes & no & greedy entropy \\
Karp \& Kleinberg~\cite{karp2007} & yes & no & noisy bisection \\
Golovin et al.~\cite{golovin2010noisy} & yes & no & EC$^2$, guaranteed \\
Kaspi et al.~\cite{kaspi2018} & yes & query-dep. & information-theoretic \\
Dawid \& Skene~\cite{dawid1979} & yes & per-observer & n/a (estimation) \\
EAR / CPR / UNICORN~\cite{lei2020ear,lei2020cpr,deng2021unicorn} & no & no & learned policy \\
Max-entropy baseline & no & no & classical IG \\
UoT~\cite{hu2024uot}, \cite{mazzaccara2024} & implicit & no & EIG over text \\
\midrule
\textbf{This work} & yes & \textbf{per-attribute} & \textbf{RAIG} \\
\bottomrule
\end{tabular}
\end{table}
